\makeatletter
\def\input@path{{./Graphics/}} % Définition du chemin d'entrée pour les graphiques
\makeatother

\documentclass[12pt,a4paper]{article} % Classe de document

% Chargement des packages essentiels
\usepackage[utf8]{inputenc} % Encodage UTF-8
\usepackage[french]{babel} % Support de la langue française
\usepackage{graphicx, amsthm, mathtools, amsfonts, amssymb, physics, longtable, mathrsfs} % Packages pour les graphiques et mathématiques

% Divers Box en couleur avec le package tcolorbox
\usepackage[many]{tcolorbox}
\usepackage{varioref}
%
% Définition avec caption et label=df:
\newtcbtheorem[number within=section]{definition}{Définition}
{fonttitle=\bfseries\upshape, fontupper=\slshape, arc=0mm, colback=blue!3,colframe=blue!75!black,separator sign dash,breakable,beamer}{df}
%
\newtcbtheorem[number within=section]{example}{Exemple}
{fonttitle=\bfseries\upshape,fontupper=\slshape, arc=0mm, colback=green!5!white,
colframe=green!55!black,separator sign dash,breakable,beamer}{exp}

\usepackage{float} % Gère le placement des objets flottants
\floatplacement{figure}{H} % Forces les figures à être placées exactement à l'endroit où elles apparaissent dans le texte

\usepackage[svgnames,x11names,rgb,table]{xcolor} % Définit les couleurs par noms
\usepackage{array, multicol, wrapfig, lmodern, multirow, booktabs, chemformula} % Diverses améliorations de mise en forme
\usepackage[detect-all=true, group-minimum-digits=4]{siunitx} % Pour les unités SI

\usepackage[bf, textfont={footnotesize,it}]{caption} % Personnalisation des légendes
\usepackage{subcaption}

% Configuration des marges de la page
\usepackage[margin=2cm, a4paper]{geometry}

\usepackage{fancyvrb} % Remplacement du mode verbatim qui permet une meilleure mise en forme
%\usepackage[Export]{adjustbox} % Utilisé pour contraindre les images à une taille maximale
%\adjustboxset{max size={0.9\linewidth}{0.9\paperheight}} % Taille maximale des images

% Configuration pour hyperref
\usepackage{hyperref} 
\hypersetup{
 colorlinks=true, % Colorise les liens
 breaklinks=true, % Permet le retour à la ligne dans les liens trop longs
 allcolors=Blue3, % Couleur des liens
 pdfauthor={Nana Engo}, % Auteur du document PDF
 pdfsubject={Chimie quantique et VQE}, % Sujet du document
 pdftitle={PHY4268 - Qubit Hamiltonian V25.04}, % Titre du document
 bookmarksopen=true % Ouvre les signets PDF par défaut
}

\usepackage{cleveref} % Améliore la gestion des références croisées

\frenchbsetup{StandardLayout=true} % Configuration du layout français
\numberwithin{equation}{section} % Numérotation des équations par section
\numberwithin{figure}{section} % Numérotation des figures par section
\numberwithin{table}{section} % Numérotation des tables par section

\setcounter{secnumdepth}{3} % Pour inclure les sous-sections
\renewcommand{\thesection}{\arabic{section}} % Numérotation des sections
%
\newcommand{\tcb}[1]{\textcolor{blue}{#1}} % textcolor{blue}{#1}
\newcommand{\tcr}[1]{\textcolor{red}{#1}} % textcolor{red}{#1}
% Titre du document
\title{\vspace*{4cm}
\Huge\tcb{Chapter 2 - Hamiltonien Qubit}
\vspace*{5cm}}

\author{
\begin{tabular}{lr}
 \textbf{Nana Engo} & \textbf{Tchapet Njafa}\\
 {\small serge.nana-engo@facsciences-uy1.cm} &\hspace*{2em}
 {\small jean-pierre.tchapet@facsciences-uy1.cm}
\end{tabular}\\[3em] % Espace entre les auteurs et la filiation
\textit{Quantum Simulation Group}\\
 \textit{Department of Physics} \\ 
 \textit{Faculty of Science} \\ 
 \textit{University of Yaoundé I}
}
\date{\vspace*{3cm}April 2025} % Date du document

\section*{Introduction}

Ce chapitre a pour objectif de présenter les bases théoriques nécessaires pour établir une correspondance entre les opérateurs fermioniques (Hamiltonien en seconde quantification) et les opérateurs qubits (Hamiltonien Qubit). Cette transformation est essentielle pour la simulation quantique des systèmes moléculaires sur des calculateurs quantiques. Nous abordons les concepts fondamentaux tels que la génération des intégrales moléculaires, la représentation nombre d'occupation, ainsi que les mappings fermion-qubit comme les transformations de Jordan-Wigner, de parité et de Bravyi-Kitaev.

À la fin de ce chapitre, l'apprenant sera capable :
\begin{itemize}
    \item de comprendre le principe d'encodage d'un problème électronique en seconde quantification sur un calculateur quantique ;
    \item de transformer un Hamiltonien moléculaire en première quantification en un Hamiltonien fermionique ;
    \item de convertir un Hamiltonien fermionique en un Hamiltonien Qubit.
\end{itemize}

\begin{figure}[h!]
    \centering
    \includegraphics[width=0.8\linewidth]{./Graphics/VQE_Flowchart_Tuto2.jpeg}
    \caption{\textit{Schéma illustrant le flux de travail du Variational Quantum Eigensolver (VQE). Les parties rouges indiquent les étapes clés liées à la transformation des Hamiltoniens.}}
    \label{fig:VQE_flowchart}
\end{figure}

\section{Hamiltonien moléculaire fermionique}

Dans cette section, nous allons explorer comment l'Hamiltonien électronique d'une molécule peut être exprimé en utilisant le formalisme de la seconde quantification. Cette approche est particulièrement utile pour les calculs de chimie quantique, car elle simplifie la manipulation des opérateurs et permet de décrire les systèmes à plusieurs électrons de manière concise. Nous allons voir comment les intégrales moléculaires, qui sont des quantités fondamentales dans les calculs de structure électronique, sont générées et utilisées pour construire l'Hamiltonien.

\subsection{Génération des intégrales moléculaires}

L'Hamiltonien électronique est exprimé dans la base des solutions de la méthode Hartree-Fock (HF), également appelées Orbitales Moléculaires (OM). Rappelons que la méthode HF est une méthode variationnelle qui permet d'approximer la fonction d'état électronique d'un système moléculaire en considérant que chaque électron se déplace dans un champ moyen créé par les autres électrons. Les orbitales moléculaires obtenues par cette méthode forment une base orthonormée qui peut être utilisée pour exprimer l'Hamiltonien électronique.

L'Hamiltonien s'écrit :
\begin{equation}\label{eq:H_elec}
\mathtt{H}_{\rm elec}=\sum_{pq} h_{pq} a^{\dagger}_p a_q + \frac12 \sum_{pqrs} h_{pqrs}  a^{\dagger}_p a^{\dagger}_q a_r  a_s,
\end{equation}
où :
\begin{itemize}
    \item $a^{\dagger}_p$ est l'\textbf{opérateur de création} d'un électron dans le spin-orbite $p$. Cet opérateur ajoute un électron à l'orbitale de spin $p$;
    \item $a_q$ est l'\textbf{opérateur d'annihilation} d'un électron dans le spin-orbite $q$. Cet opérateur supprime un électron de l'orbitale de spin $q$;
    \item $h_{pq}$ sont les \textbf{intégrales à un électron}, qui décrivent l'énergie cinétique de l'électron et son interaction avec le potentiel créé par les noyaux atomiques;
    \item $h_{pqrs}$ sont les \textbf{intégrales à deux électrons}, qui décrivent l'interaction coulombienne entre les électrons.
\end{itemize}

Les intégrales à un électron sont données par :
\begin{equation}\label{eq:h_pq}
h_{pq} = \int \phi^*_p(\mathbf{r}) \left( -\frac{1}{2} \nabla^2 - \sum_{I} \frac{Z_I}{|\mathbf{r}-\mathbf{R}_I|} \right)   \phi_q(\mathbf{r})d\mathbf{r},
\end{equation}
où :
\begin{itemize}
    \item $\phi_p(\mathbf{r})$ est l'orbitale moléculaire $p$ au point $\mathbf{r}$;
    \item $Z_I$ est la charge du noyau $I$;
    \item $\mathbf{R}_I$ est la position du noyau $I$.
\end{itemize}
Ces intégrales décrivent l’énergie cinétique des électrons individuels et leurs interactions Coulombiennes avec les champs électriques du noyau.

Les intégrales à deux électrons sont données par :
\begin{equation}\label{eq:h_pqrs}
h_{pqrs} = \int \frac{\phi^*_p(\mathbf{r}_1)  \phi^*_q(\mathbf{r}_2) \phi_r(\mathbf{r}_2)  \phi_s(\mathbf{r}_1)}{|\mathbf{r}_1-\mathbf{r}_2|}d\mathbf{r}_1d\mathbf{r}_2,
\end{equation}
où $\mathbf{r}_1$ et $\mathbf{r}_2$ sont les positions des deux électrons. Ces intégrales décrivent les interactions Coulombiennes entre les électrons.

Il est important de noter que l'opérateur $a^{\dagger}_p a_q $ est l'\textbf{opérateur d'excitation}, qui excite un électron du spin-orbite occupé $\phi_q$ vers le spin-orbite inoccupé $\phi_p$.

\bigskip

\subsection{Orbitales moléculaires et orbitales de spin} 

Les orbitales moléculaires ($\phi_u$) peuvent être occupées ou virtuelles (inoccupées). Un OM peut contenir au maximum deux électrons (avec des spins opposés) conformément au principe d'exclusion de Pauli. Ces OM sont les solutions de la méthode HF, où chaque électron est soumis au champ moyen créé par les autres électrons.

Cependant, dans ce qui suit, nous travaillons en fait avec des \textbf{orbitales de spin}, qui sont associées à un spin up ($\alpha$) ou spin down ($\beta$). Ces deux spins sont également communément désignés par $\alpha$ et $\beta$, respectivement. Ainsi, les orbitales de spin peuvent contenir un électron ou être inoccupées.

Le \Cref{tab:1_2_Quant} résume les correspondances entre les éléments du Hamiltonien moléculaire dans les formalismes de première et seconde quantification.

\begin{table}
\begin{tabular}{|l|l|l|}\hline
\textbf{Approximation BO}  &  \textbf{Combinaison en 2e quantification} &  \textbf{Description} \\ \hline
$-\frac12\sum_i\nabla^2_i$ & $-\frac12\sum_{p,q}\langle p|\nabla^2_p|q\rangle a^\dagger_p a_q$ & $\mathtt{T}_e$ (énergie cinétique) \\ \hline
$-\sum_{i,I}\frac{Z_I}{|\mathbf{r}_i-\mathbf{R}_I|}$ &
$-\sum_{p,q} \langle p|\frac{Z_I}{|\mathbf{r}_p-\mathbf{R}_I|}|q\rangle a^\dagger_p a_q$ & $\mathtt{U}_{en}$ (interaction noyau-électron) \\ \hline
$\sum_{i,j>i}\frac{1}{|\mathbf{r}_i-\mathbf{r}_j|}$ & 
$\sum_{p,q,r,s}\langle pq \big|\frac{1}{|\mathbf{r}_i-\mathbf{r}_j|}\big|rs\rangle a^\dagger_p a^\dagger_q a_r a_s$ & 
$\mathtt{U}_{ee}$ (interaction électron-électron) \\ \hline
\end{tabular}
\caption{Correspondance Première et Seconde Quantification}
\label{tab:1_2_Quant}
\end{table}

Un des avantages importants apporté par le formalisme de seconde quantification est que la propriété d'anti-symétrie des fonctions d'état pour l'échange de fermions identiques, prise en compte manuellement dans une approche de première quantification en considérant les déterminants de Slater, est automatiquement imposée par les relations d'anti-commutation des opérateurs de création et d'annihilation. Cependant, le prix à payer est de travailler avec un nombre non fixe de particules, qui reste de toute façon une quantité conservée.

\subsection{Répresentation nombre d'occupation}

La représentation nombre d'occupation est une manière concise de décrire l'état d'un système de fermions (comme les électrons) en spécifiant le nombre de particules présentes dans chaque orbitale de spin. Au lieu de suivre chaque électron individuellement, on se concentre sur le nombre d'électrons dans chaque état possible. Cette représentation est particulièrement utile en seconde quantification, car elle permet de manipuler facilement les opérateurs de création et d'annihilation.

En seconde quantification, le déterminant de Slater est représenté par un \textbf{vecteur nombre d'occupation}:
\begin{equation}\label{eq:occupation_vector}
\ket{f} = \ket{f_{M-1}, \dots, f_{i-1}, f_i, f_{i+1}, \dots, f_0},
\end{equation}
où :
\begin{itemize}
    \item $f_i = 1$ si le spin-orbite $\phi_i$ est occupé (c'est-à-dire qu'il contient un électron) et donc présent dans le déterminant de Slater.
    \item $f_i = 0$ si le spin-orbite $\phi_i$ est vide (c'est-à-dire qu'il ne contient pas d'électron) et donc absent du déterminant de Slater.
\end{itemize}

Le vecteur $\ket{f_i}$ est appelé \textbf{vecteur nombre d'occupation de l'orbitale fermionique i}, parce que :
\begin{align}
\label{eq:occupation_states}
& \ket{0} = \begin{pmatrix} 1 \\ 0 \end{pmatrix} : \text{état vide}, \\
& \ket{1} = \begin{pmatrix} 0 \\ 1 \end{pmatrix} : \text{état occupé}.
\end{align}

Dans la deuxième représentation de quantification, au lieu de poser la question \textit{Quel électron est dans quel état?}, nous posons la question \textit{combien de particules y a-t-il dans chaque état?}.

\subsection{Opérateurs fermioniques}

Les opérateurs fermioniques, en particulier les opérateurs de création ($a^\dagger$) et d'annihilation ($a$), sont des outils mathématiques essentiels pour décrire les systèmes de fermions en seconde quantification. Ils permettent de manipuler les états quantiques en ajoutant ou en supprimant des particules, tout en respectant l'antisymétrie requise par la nature fermionique des électrons. Dans cette section, nous allons définir ces opérateurs et examiner leurs propriétés fondamentales.

Les opérateurs fermioniques $a^\dagger_i$ et $a_i$ obéissent aux relations d'anticommutation fermioniques suivantes :
\begin{align}
\label{eq:anticommutation}
& \{a_p, a^\dagger_q\} = a_p a^\dagger_q + a^\dagger_q a_p = \delta_{pq}, \\
& \{a_p, a_q\} = \{a^\dagger_p, a^\dagger_q\} = 0,
\end{align}
où $\delta_{pq}$ est le symbole de Kronecker (égal à 1 si $p = q$ et 0 sinon).

La nature fermionique des électrons implique que les fonctions d'état à plusieurs électrons doivent être antisymétriques par rapport à l'échange de particules. Cela se reflète dans la manière dont les opérateurs de création et d'annihilation agissent sur les déterminants $\ket{f}$ :
\begin{align}
\label{eq:fermion_operators}
& a_i \ket{f_{M-1}, \dots, f_{i-1}, f_i, f_{i+1}, \dots, f_0} = \delta_{f_i, 1} (-1)^{\sum_{m < i} f_m} \ket{f_{M-1}, \dots, f_{i-1}, f_i \oplus 1, f_{i+1}, \dots, f_0}, \\
& a^\dagger_i \ket{f_{M-1}, \dots, f_{i-1}, f_i, f_{i+1}, \dots, f_0} = \delta_{f_i, 0} (-1)^{\sum_{m < i} f_m} \ket{f_{M-1}, \dots, f_{i-1}, f_i \oplus 1, f_{i+1}, \dots, f_0}, \\
& a \ket{0} = a^\dagger \ket{1} = 0.
\end{align}
où :
\begin{itemize}
    \item Le terme de phase $p_i = (-1)^{\sum_{m < i} f_m}$ désigne la parité qu'impose l'anti-symétrie d'échange des fermions.
        \begin{itemize}
            \item $p_i = -1$ si le nombre d'électrons est impair dans les orbitales de spin d'indice inférieur à $i$.
            \item $p_i = 1$ si le nombre d'électrons est pair dans les orbitales de spin d'indice inférieur à $i$.
        \end{itemize}
    \item Le symbole $\oplus$ représente l'addition modulo 2, c'est-à-dire $1 \oplus 1 = 0$ et $0 \oplus 1 = 1$.
    \item L'opérateur de création $a^\dagger_i$ ajoute un électron dans une orbitale inoccupée $i$ (ou dans $\ket{f_i}$).
    \item L'opérateur d'annihilation $a_i$ supprime un électron dans une orbitale occupée $i$ (ou en $\ket{f_i}$).
    \item $a_i^2 = (a^\dagger_i)^2 = 0$ pour tout $i$. On ne peut pas créer ou anéantir un fermion dans le même mode deux fois.
\end{itemize}

Par exemple :
\begin{align}
\label{eq:fermion_example}
& a^\dagger_1 \ket{0}_1 = \ket{1}_1, \\
& a^\dagger_1 \ket{1}_1 = (a^\dagger_i)^2 \ket{0}_1 = 0, \\
& a_1 \ket{1}_1 = \ket{0}_1, \\
& a_1 \ket{0}_1 = a_1^2 \ket{1}_1 = 0.
\end{align}
Notez qu'ici $a^\dagger_1 \ket{1}_1 = 0$ et $a_1 \ket{0}_1 = 0$ signifient le vecteur nul (zéro) et non $\ket{0}_1$.

En utilisant également de tels opérateurs, nous pouvons exprimer :
\begin{equation}
\label{eq:example_expression}
\ket{0} \ket{1} \ket{1} \ket{0} = a^\dagger_1 a^\dagger_2 \ket{0}^{\otimes 4}.
\end{equation}

L'opérateur d'occupation orbitale est donné par :
\begin{align}
\label{eq:occupation_operator}
& n_i = a^\dagger_i a_i, \\
& n_i \ket{f_{M-1}, \dots, f_i, \dots, f_0} = f_i \ket{f_{M-1}, \dots, f_i, \dots, f_0},
\end{align}
et il compte le nombre d'électrons dans une orbitale donnée, c'est-à-dire :
\begin{align}
\label{eq:occupation_number}
& n_i \ket{0}_i = 0, \\
& n_i \ket{1}_i = \ket{1}_i.
\end{align}

\subsection{Exercise}

\textbf{Enoncé :} \\
Dans l'\textbf{Exemple d'illustration} ci-dessus, nous avons considéré un système fictif décrit par des orbitales de spin $\ket{A_\uparrow}$, $\ket{A_\downarrow}$, $\ket{B_\uparrow}$, $\ket{B_\downarrow}$ étiquetées comme suit :
\begin{align}
\label{eq:spin_orbitals}
& \ket{A_\uparrow} = \ket{00} = \ket{\mathbf{0}}, \\
& \ket{A_\downarrow} = \ket{01} = \ket{\mathbf{1}}, \\
& \ket{B_\uparrow} = \ket{10} = \ket{\mathbf{2}}, \\
& \ket{B_\downarrow} = \ket{11} = \ket{\mathbf{3}}.
\end{align}
Nous supposons que l'état Hartree-Fock pour ce système \textbf{a les deux électrons occupant les orbitales $\ket{A}$}. Nous stockons les occupations des orbitales de spin $\ket{A_\uparrow}$, $\ket{A_\downarrow}$, $\ket{B_\uparrow}$, $\ket{B_\downarrow}$ que nous ordonnons comme $\ket{f_{B_\downarrow}, f_{B_\uparrow}, f_{A_\downarrow}, f_{A_\uparrow}}$, avec $f_i = 0, 1$ (dans l'espace de Fock fermionique).

\begin{enumerate}
    \item Écrivez, dans l'espace de Fock fermionique, l'état Hartree-Fock $\ket{\Psi_{\rm HF}}$, c'est-à-dire la correspondance du déterminant de Slater antisymétrisé obtenu dans l'\textbf{Exemple d'illustration}.
    \begin{equation}
    \label{eq:hf_state}
    \ket{\Psi_{\rm HF}} = \ket{0011}.
    \end{equation}
    \item Écrivez, dans l'espace de Fock fermionique, la fonction d'état $s_z = 0$.
    \begin{equation}
    \label{eq:sz_state}
    \ket{\Psi_{\rm HF}} = \alpha \ket{0011} + \beta \ket{1100} + \ket{1001} + \ket{0101}.
    \end{equation}
\end{enumerate}

\section{Hamiltonien Qubit}

\subsection{Correspondance (mapping) fermion-qubit}

\noindent \textbf{Introduction :} \\
Pour simuler des systèmes quantiques sur des ordinateurs quantiques, il est nécessaire de transformer les opérateurs fermioniques, qui décrivent les électrons dans les molécules, en opérateurs qubit, qui peuvent être implémentés sur les qubits des ordinateurs quantiques. Cette transformation est connue sous le nom de "mapping fermion-qubit". Dans cette section, nous allons explorer les principales méthodes de mapping fermion-qubit, telles que les transformations de Jordan-Wigner, de parité et de Bravyi-Kitaev.

\bigskip

En simulation quantique, encoder un problème de structure électronique en seconde quantification sur un calculateur quantique revient à établir une correspondance entre les opérateurs d'échelle fermioniques et les opérateurs de Pauli. Les transformations de Jordan-Wigner (JWT), de Parité (PT) et de Bravyi-Kitaev (BKT) sont parmi les plus utilisées en calculs quantiques. Pour plus de détails, consulter J. T. Seeley, M. J. Richard, and P. J. Love, "\textit{The Bravyi-Kitaev transformation for quantum computation of electronic structure}", J. Chem. Phys. \textbf{137}, 224109 (2012).

\subsubsection{Décomposition de Pauli}

\noindent \textbf{Introduction :} \\
Les matrices de Pauli sont un ensemble de matrices hermitiennes et unitaires qui forment une base complète pour l'espace des opérateurs agissant sur un qubit. Cela signifie que tout opérateur agissant sur un qubit peut être exprimé comme une combinaison linéaire des matrices de Pauli. Cette propriété est cruciale pour la simulation quantique, car elle permet de décomposer l'Hamiltonien électronique en termes d'opérateurs simples qui peuvent être implémentés sur les qubits.

\bigskip

Comme $\{\mathbb{I}, \mathtt{X}, \mathtt{Y}, \mathtt{Z}\}$ forme une base complète pour tout opérateur Hermitien 1- ou multi-qubit, l'ingrédient clé des algorithmes variationnels est de décomposer l'Hamiltonien électronique en termes de produits des matrices de Pauli :
\begin{equation}
\label{eq:pauli_decomposition}
\mathtt{P}_j = \mathtt{X}_{M-1}^j \otimes \mathtt{X}_{M-2}^j \otimes \mathtt{X}_k^j \otimes \dots \otimes \mathtt{X}_0^j = \bigotimes_{i=0}^{M-1} \mathtt{X}_i^j,
\end{equation}
où les opérateurs de Pauli $\mathtt{X}^j \in \{\mathbb{I}, \mathtt{X}, \mathtt{Y}, \mathtt{Z}\}$ sont tels que :
\begin{equation}
\label{eq:pauli_algebra}
\begin{aligned}
& \mathtt{X}^i \mathtt{X}^j = \mathbb{I} \delta_{ij} + i \varepsilon_{ijk} \mathtt{X}^k, \\
& \varepsilon_{ijk} = \begin{cases}
+1, & \text{pour les permutations circulaires droites de } (i, j, k) \\
-1, & \text{pour les permutations circulaires gauches de } (i, j, k) \\
0, & \text{sinon}
\end{cases}
\end{aligned}
\end{equation}

\begin{align}
\label{eq:pauli_matrices}
& \mathtt{X} := \ket{0}\bra{1} + \ket{1}\bra{0}, \\
& i \mathtt{Y} := \ket{0}\bra{1} - \ket{1}\bra{0}, \\
& \mathtt{Z} := \ket{0}\bra{0} - \ket{1}\bra{1}.
\end{align}
Les opérateurs de Pauli sont à la fois unitaires et Hermitiens. On a utilisé la notation de Dirac pour les opérateurs.
On note également que :
\begin{align}
\label{eq:pauli_matrices_simplifiees}
& \mathtt{X} := \kb{0}{1} + \kb{1}{0}, \\
& i \mathtt{Y} := \kb{0}{1} - \kb{1}{0}, \\
& \mathtt{Z} := \proj{0} - \proj{1}.
\end{align}

L'Hamiltonien total peut alors être représenté comme la combinaison linéaire de $\mathtt{P}_j$ :
\begin{align}
\label{eq:hamiltonian_pauli}
& \mathtt{H} = \sum_{j=1}^{N_k} w_j \mathtt{P}_j = \sum_{j=1}^{N_k} w_j \left( \bigotimes_{i=0}^{M-1} \mathtt{X}_i^j \right), \\
& w_j = \langle \mathtt{H}, \mathtt{P}_j \rangle = \text{Tr}(\mathtt{H}^\dagger \mathtt{P}_j),
\end{align}
où $w_j$ sont les poids de la chaîne de Pauli $\mathtt{P}_j$, $i$ indique sur quel qubit l'opérateur agit, et $j$ désigne le terme dans l'Hamiltonien.

Pour un Hamiltonien de la 2e quantification, le nombre total de chaînes de Pauli $N_k$ dépend du nombre de termes $h_{pq}$ à 1-électron et $h_{pqrs}$ à 2-électrons.

\subsubsection{Opérateurs d'échelle qubit}

Les opérateurs d'échelle qubit sont des opérateurs qui agissent localement sur les qubits et qui permettent de créer ou d'annihiler des excitations dans le système quantique. Ils sont analogues aux opérateurs de création et d'annihilation fermioniques, mais agissent sur les qubits au lieu des orbitales de spin. Dans cette section, nous allons définir ces opérateurs et examiner leurs propriétés.

\bigskip

Ce sont les opérateurs suivants qui agissent localement sur les qubits :
\begin{equation}
\label{eq:qubit_operators}
\begin{array}{l|l}
\textbf{Opérateur qubit} & \textbf{Description} \\ \hline
\mathbb{I} := \ket{0}\bra{0} + \ket{1}\bra{1} = \begin{pmatrix} 1 & 0 \\ 0 & 1 \end{pmatrix} & \text{Identité} \\
\mathtt{Q}^- := \ket{0}\bra{1} = \begin{pmatrix} 0 & 1 \\ 0 & 0 \end{pmatrix} = \frac{1}{2}(\mathtt{X} + i\mathtt{Y}) & \text{Annihilation} \\
\mathtt{Q}^+ := \ket{1}\bra{0} = \begin{pmatrix} 0 & 0 \\ 1 & 0 \end{pmatrix} = \frac{1}{2}(\mathtt{X} - i\mathtt{Y}) & \text{Création} \\
\mathtt{Q}^+ \mathtt{Q}^- := \ket{1}\bra{1} = \begin{pmatrix} 0 & 0 \\ 0 & 1 \end{pmatrix} = \frac{1}{2}(\mathbb{I} - \mathtt{Z}) & \text{Un nombre (particule)} \\
\mathtt{Q}^- \mathtt{Q}^+ := \ket{0}\bra{0} = \begin{pmatrix} 1 & 0 \\ 0 & 0 \end{pmatrix} = \frac{1}{2}(\mathbb{I} + \mathtt{Z}) & \text{Zéro nombre (trou)} \\ \hline
\end{array}
\end{equation}
On a utilisé la notation de Dirac pour les opérateurs. On note également que :
\begin{equation}
\label{eq:qubit_operators_simplifies}
\begin{array}{l|l}
\textbf{Opérateur qubit} & \textbf{Description} \\ \hline
\mathbb{I} := \proj{0} + \proj{1} = \begin{pmatrix} 1 & 0 \\ 0 & 1 \end{pmatrix} & \text{Identité} \\
\mathtt{Q}^- := \kb{0}{1} = \begin{pmatrix} 0 & 1 \\ 0 & 0 \end{pmatrix} = \frac{1}{2}(\mathtt{X} + i\mathtt{Y}) & \text{Annihilation} \\
\mathtt{Q}^+ := \kb{1}{0} = \begin{pmatrix} 0 & 0 \\ 1 & 0 \end{pmatrix} = \frac{1}{2}(\mathtt{X} - i\mathtt{Y}) & \text{Création} \\
\mathtt{Q}^+ \mathtt{Q}^- := \proj{1} = \begin{pmatrix} 0 & 0 \\ 0 & 1 \end{pmatrix} = \frac{1}{2}(\mathbb{I} - \mathtt{Z}) & \text{Un nombre (particule)} \\
\mathtt{Q}^- \mathtt{Q}^+ := \proj{0} = \begin{pmatrix} 1 & 0 \\ 0 & 0 \end{pmatrix} = \frac{1}{2}(\mathbb{I} + \mathtt{Z}) & \text{Zéro nombre (trou)} \\ \hline
\end{array}
\end{equation}

Les opérateurs qubits sont antisymétriques : $\{\mathtt{Q}^+, \mathtt{Q}^-\} = \mathtt{Q}^+ \mathtt{Q}^- + \mathtt{Q}^- \mathtt{Q}^+ = 1$ pour la même spin-orbitale, de façon similaire aux opérateurs fermioniques.

\section{Hamiltonien Qubit}

Pour simuler l'Hamiltonien moléculaire sur un calculateur quantique, il est nécessaire de convertir les opérateurs fermioniques en opérateurs qubits. Cette conversion peut être réalisée à l'aide des transformations suivantes :
\begin{itemize}
    \item La transformation de Jordan-Wigner (JWT) ;
    \item La transformation de parité (PT) ;
    \item La transformation de Bravyi-Kitaev (BKT).
\end{itemize}

La décomposition en termes d'opérateurs de Pauli est essentielle dans ce contexte. Par exemple, tout opérateur Hermitien peut être exprimé comme une combinaison linéaire des matrices $\{\mathbb{I}, X, Y, Z\}$ :
\begin{equation}
H = \sum_j w_j P_j,
\end{equation}
où $P_j$ représente un produit tensoriel d'opérateurs de Pauli agissant sur différents qubits.

Considérons un exemple simple où l'Hamiltonien fermionique est transformé en termes d'opérateurs qubits via la transformation JWT. Les détails mathématiques sont laissés en exercice pour l'apprenant.

\subsection{Correspondance fermion-qubit}\label{sec:mapping}

Les trois principales transformations pour encoder les opérateurs fermioniques en qubits sont :

\begin{itemize}
    \item \textbf{Transformation de Jordan-Wigner (JWT)} : 
    \begin{equation}
        a_j^\dagger \mapsto \frac{1}{2} \left(\bigotimes_{k=0}^{j-1} Z_k \right) \otimes (X_j - iY_j)
    \end{equation}
    
    \item \textbf{Transformation de parité (PT)} :
    \begin{equation}
        a_j^\dagger \mapsto \frac{1}{2} \left(\bigotimes_{k=0}^{j-1} X_k \right) \otimes (X_j - iY_j) \otimes \left(\bigotimes_{k=j+1}^{N-1} Z_k \right)
    \end{equation}
    
    \item \textbf{Transformation de Bravyi-Kitaev (BKT)} :
    \begin{equation}
        a_j^\dagger \mapsto \frac{1}{2} \left(\bigotimes_{k \in \partial(j)} X_k \right) \otimes (X_j - iY_j)
    \end{equation}
\end{itemize}

\begin{figure}[h!]
    \centering
    \includegraphics[width=0.7\linewidth]{./Graphics/mapping_comparison.png}
    \caption{\textit{Comparaison des différentes transformations fermion-qubit pour un système à 4 qubits. Les couleurs représentent les types d'opérateurs Pauli impliqués.}}
    \label{fig:mapping_compare}
\end{figure}

\subsection{Exemple : Molécule d'hydrogène}\label{sec:H2_example}

Considérons la molécure H$_2$ dans une base STO-3G. L'Hamiltonien fermionique se réduit à :

\begin{equation}
H = -1.05\ \mathbb{I} + 0.79\ Z_0 + 0.79\ Z_1 + 0.60\ Z_0Z_1 + 0.20\ Y_0Y_1
\end{equation}

\begin{table}[h!]
    \centering
    \caption{Comparaison des mappings pour H$_2$ (6 qubits)}
    \begin{tabular}{lccc}
        \toprule
        \textbf{Métrique} & JWT & PT & BKT \\
        \midrule
        Nombre de termes & 14 & 11 & 9 \\
        Profondeur CNOT & 56 & 48 & 32 \\
        Fidélité & 0.992 & 0.994 & 0.997 \\
        \bottomrule
    \end{tabular}
    \label{tab:H2_comparison}
\end{table}

% Section : Implémentation quantique
\section{Implémentation sur calculateur quantique}\label{sec:implementation}

Le circuit quantique typique pour la simulation VQE comprend trois étapes clés :

\begin{enumerate}
    \item \textbf{Encodage de l'état} : Préparation de l'état Hartree-Fock
    \item \textbf{Ansatz} : Circuit paramétrique (ex: UCCSD)
    \item \textbf{Mesure} : Estimation des valeurs moyennes des termes de Pauli
\end{enumerate}

\begin{figure}[h!]
    \centering
    \includegraphics[width=0.9\linewidth]{./Graphics/VQE_circuit.png}
    \caption{\textit{Circuit quantique complet pour la simulation VQE de H$_2$. Les parties en rouge montrent l'ansatz UCCSD, et les mesures en bleu correspondent aux termes de Pauli.}}
    \label{fig:VQE_circuit}
\end{figure}


\section*{Conclusion}

Ce chapitre a présenté les fondements théoriques et pratiques de la conversion des Hamiltoniens moléculaires en formalismes adaptés aux calculateurs quantiques. Les points clés à retenir sont :

\begin{itemize}
    \item La seconde quantification permet une description compacte des systèmes électroniques
    \item Les mappings fermion-qubit préservent les relations d'anticommutation
    \item Le choix du mapping influence directement les performances du calcul quantique
\end{itemize}

Ces concepts seront mis en pratique dans le prochain chapitre à travers l'implémentation complète d'un algorithme VQE pour la molécure de H$_2$ sur un simulateur quantique.

% Exercices
\section*{Exercices}

\begin{example}{Application de la JWT}{ex_jwt}
Convertir l'opérateur fermionique $a_2^\dagger a

% Conclusion
\section*{Conclusion}

Dans ce chapitre, nous avons exploré les bases théoriques nécessaires pour établir une correspondance entre les opérateurs fermioniques et qubits. Cette étape est cruciale pour la simulation quantique des systèmes moléculaires. Les transformations telles que Jordan-Wigner, parité et Bravyi-Kitaev permettent cette conversion tout en préservant les propriétés fondamentales du système.

Les concepts abordés ici serviront de fondement pour comprendre et implémenter des algorithmes quantiques comme le Variational Quantum Eigensolver (VQE). Dans le prochain chapitre, nous approfondirons l'application pratique de ces concepts dans le cadre des simulations moléculaires.

\vspace*{\fill}

% Fin du document
\end{document}

